\section{Introduction}\label{sec:intro}

Short-term rental platforms have fundamentally reshaped urban housing markets over the past decade. By allowing property owners to switch flexibly between short-term and long-term rental uses, platforms such as Airbnb have introduced a new form of market arbitrage: units flow to the highest-yielding use in each period, which in high-demand tourist and business travel markets is often a nightly rental rather than a monthly lease. A substantial empirical literature has documented that this reallocation reduces the supply of long-term rental housing and increases market rents in affected neighborhoods. Yet the policy response has been uneven: cities have adopted a patchwork of regulations ranging from platform bans and licensing caps to neighborhood-level registration requirements, and rigorous causal evidence on how specific regulatory instruments affect housing outcomes remains thin.

This paper studies one such instrument --- Chicago's Shared Housing Ordinance prohibited buildings registry --- which creates a building-level, permanent opt-out mechanism for STR activity. Unlike city-wide platform bans, which affect all neighborhoods simultaneously and confound identification with macro shocks, Chicago's mechanism generates staggered variation: individual buildings adopt prohibitions at different times from 2015 through the present, and the resulting variation in census tract-level prohibition intensity evolves gradually over a decade. This variation allows us to apply modern difference-in-differences methods that explicitly account for treatment effect heterogeneity across cohorts, delivering more credible causal estimates than the standard two-way fixed effects regression.

\paragraph{Literature on STR platforms and housing markets.} A growing body of research documents that short-term rental platforms exert upward pressure on rents through a supply reduction channel. \citet{barron2021effect} exploit Airbnb's staggered market entry across U.S. cities and find that a 1 percent increase in Airbnb listings raises local rents by 0.018 percent and house prices by 0.026 percent, with effects concentrated in neighborhoods close to tourist attractions. \citet{koster2021too} use the abrupt entry of Airbnb into the Amsterdam market to estimate that high-density listing areas experience rent increases of 4--6 percent relative to less-affected areas. \citet{horn2017is} find similar effects in Boston, with stronger impacts in ZIP codes with above-median listing growth. Together, these studies establish a plausible mechanism: STR platforms convert long-term rental units to short-term use, tightening the local rental market and driving up rents for conventional tenants.

The natural policy implication is that restricting STR activity should relieve rental supply pressure and reduce rents. Several papers examine this prediction for city-level regulations. \citet{garcia-lopez2020airbnb} study Barcelona's 2014 and 2016 zoning-based STR restrictions and find that affected areas see rent reductions of 4--7 percent over the subsequent three years. However, city-wide restrictions applied simultaneously to all neighborhoods in a city offer limited identification variation beyond the regulated-vs.-unregulated city comparison, which is vulnerable to the same unobserved urban trends that afflict non-experimental comparisons in general.

The evidence on more targeted, neighborhood-level restriction mechanisms is much thinner. Studies of Chicago's ordinance specifically are absent from the peer-reviewed literature, despite Chicago operating one of the largest and most granular prohibited buildings datasets in the country. Our paper fills this gap.

\paragraph{Econometric methodology.} The identification challenge in staggered adoption settings is well understood in the recent econometrics literature. \citet{goodman-bacon2021difference} shows that the two-way fixed effects estimator decomposes into a weighted sum of all pairwise two-by-two comparisons, with some weights potentially negative when treatment effects vary across cohorts. \citet{de-chaisemartin2020two} provide complementary conditions under which TWFE estimates can have the wrong sign even if every individual treatment effect is positive. These concerns are directly relevant to our setting, where early-adopting tracts in gentrifying neighborhoods may respond differently to prohibitions than late-adopting tracts.

We address these concerns using the Callaway--Sant'Anna (CS) estimator \citep{callaway2021difference}, which constructs group-time average treatment effects using only valid comparison units --- never-treated tracts in our preferred specification --- and aggregates them with transparent, non-negative weights. We supplement this with the sensitivity framework of \citet{rbs2023}, which bounds the identified set for the ATT under departures from parallel trends parameterized by a smoothness constraint. Together, these tools allow us to be explicit about the degree to which our findings depend on assumptions about counterfactual rent paths.

\paragraph{This paper's contribution.} We make three contributions. First, we provide the first causal estimates of the rent effect of Chicago's building-level STR prohibition mechanism, exploiting ten years of staggered prohibition adoption across 842 census tracts. Second, we document a V-shaped pre-treatment divergence between treated and never-treated tracts that reveals the endogenous targeting of prohibitions: prohibition activity concentrated in tracts that were already on faster rent-growth trajectories, consistent with high-demand neighborhoods being most motivated to restrict STR activity. This selection finding is itself policy-relevant, as it cautions against interpreting any rent difference between prohibited and unrestricted tracts as a causal effect without accounting for pre-existing trajectories. Third, we apply HonestDiD sensitivity analysis to show that our results remain consistent with positive rent effects even under the most pessimistic parallel trends violation that is consistent with the observed pre-period dynamics.

\paragraph{Main findings.} Our baseline two-feature matched CS ATT is $+\$4.60$/month (SE 0.39); the five-feature strict matched ATT is $+\$5.46$/month; the full-panel CS estimate is $+\$20.65$/month. The binary-treatment sensitivity specification yields $+\$6.4$/month, close to the published clinic benchmark. Neither estimate should be taken as a precise causal quantity given the parallel trends violation; all are consistent with positive effects under HonestDiD sensitivity analysis with smoothness bound $M = 2.71$. Design-parameter robustness grids across matching depth ($k \in \{1, \ldots, 5\}$) and treatment intensity thresholds ($p \in \{0.10, \ldots, 0.50\}$) show stable positive signs throughout. A SUTVA diagnostic using the donut design finds that excluding never-treated tracts that are adjacent to treated tracts raises the estimated ATT substantially (to $+\$26$/month for isolated controls), consistent with positive spillovers that attenuate the headline estimate downward.

\paragraph{Paper organization.} Section~\ref{sec:policy} describes Chicago's Shared Housing Ordinance and the prohibited buildings mechanism. Section~\ref{sec:data} presents the data sources, spatial crosswalk, and panel construction. Section~\ref{sec:methods} develops the empirical strategy, including the CS estimator, the parallel trends challenge, and the matching procedure. Section~\ref{sec:results} reports the results across all specification layers. Section~\ref{sec:limits} discusses threats to validity and limitations. Section~\ref{sec:conclusion} concludes with policy implications and directions for future research.
